\begin{lemma}[Scheide structural reduction and exit-color bijection (1)]
Every Scheide structural configuration admits a reordered tree-sequence
$T'$ and reduced configuration on $T'$ with $V(T')=V(T)$.  Every old base
vertex remains in the new base.  For each old block, either its color,
two vertex identities, and two edge identities occur together as a new
block, or both of its vertices have been absorbed into the new base.
Consequently the exits are in bijection with the remaining, pairwise
distinct block colors as specified by the reduced configuration.
\end{lemma}
\begin{proof}
Put $X=V(T)$.  By the globally maximum clause of
\noderef{ScheideStructuralConfiguration} and
\noderef{ScheideGlobalMaximumOrder}, $X$ is elementary.  Hence each block
color is missing at exactly one vertex of $X$: existence is condition (S2),
and two such vertices would violate elementariness.

For every exit vertex $v$, condition (S4) in
\noderef{ScheideStructuralConfiguration} supplies at least one block color
missing at $v$; choose one and call it $c(v)$.  If $v\ne w$, then
$c(v)\ne c(w)$, again by elementariness.  Thus the selected colors form a
set $\Gamma'$ of cardinality equal to the number of exits, and every exit
misses exactly one color of $\Gamma'$.

Retain precisely the blocks whose colors lie in $\Gamma'$.  Move all other
blocks before them and regard their two vertices, in their displayed order,
as the end of the enlarged base segment.  This is still a Tashkinov
tree-sequence: each of the two edges of a moved block joins its new vertex
to the old base by (S1), and its common color is missing on the old base by
(S2), so both attachment and missing-color witnesses remain valid.  The
operation only permutes named block pairs and changes where the base ends;
it does not change the coloring or any vertex or edge identity.  Therefore
the new sequence $T'$ has exactly $V(T)$, contains every old base vertex in
its new base, and each old block is either retained with all four names or
has both vertices in that base.

Conditions (S1), (S2), and (S4) are unchanged for retained blocks.  For
(S3), the newly incorporated block edges have colors in
$\Gamma\setminus\Gamma'$, hence outside the new block-color set
$\Gamma'$.  Thus all fields of
\noderef{ScheideStructuralConfiguration} hold for the reordered data.
The injection $v\mapsto c(v)$ has domain and codomain of the same finite
cardinality, so it is a bijection.  Surjectivity says every retained block
color is missing at an exit; injectivity plus elementariness gives the
unique-color clause at every exit.  These are exactly the additional fields
of \noderef{ScheideReducedConfiguration}, proving the asserted existence
and all identity-preservation clauses.
\end{proof}
